\scanpage{183}
\mspage{132}
On a alors le

\uline{Corollaire 1} \inserted{Supposons $W(7\,\mathrm{bis})$.} Si le \deleted{diagramme} \inserted{carré} $(\ast)$ est \inserted{$W$-}homotopiquement cocartésien, alors
\[
\left\{\begin{aligned}
\alpha_1\in W&\Longleftrightarrow f_2\in W\\
\alpha_2\in W&\Longleftrightarrow f_1\in W .
\end{aligned}\right.
\]
Cela résulte aussitôt du diagramme cubique de la page précédente, \deleted{\unsure{où}} \deleted{\unsure{toutes}} les flèches verticales sont dans $W$, et de l'énoncé similaire concernant la face supérieure du cube, donné dans la prop.~13 p.~125.

\uline{Corollaire 2} \inserted{Supposons $W(7)$.} Considérons un diagramme cubique
\diagram{\begin{tikzcd}[ampersand replacement=\&,column sep=small,row sep=small]
X_0\arrow[r,"\alpha_1"]\arrow[rrd,"\alpha_2"]\arrow[ddd,-,"\deleted{f}\unsure{p|\rho|\varphi}_0"']\&X_1\arrow[rrd,"f_1"]\arrow[ddd,"\deleted{f}p_1",crossing over]\&\&\\
\&\&X_2\arrow[r,-,"f_2"]\arrow[ddd,"\deleted{f}p_2",crossing over]\&Y\arrow[ddd,"\deleted{f}p"]\\
\&\&\&\\
X'_0\arrow[r,-,"\alpha'_1"]\arrow[rrd,-,"\alpha'_2"]\&X'_1\arrow[rrd,"f'_1"]\&\&\\
\&\&X'_2\arrow[r,"f'_2"']\&Y'\rlap{,}
\end{tikzcd}}
\deleted{\unsure{où}} où les deux faces inf. et sup. du cube sont des carrés $W$-homotopiquement cocartésiens. Considérons les conditions
% archive sheet number: 183
